\scanpage{14}

$C\in\Delta_\sigma$ ssi

a) si $\unsure{\varphi}\ (R,\delta\,;\bar R,\bar\delta,\varphi:R\unsure{-|\to}\bar R)\in\sigma_1$, alors
\[
\mathrm{Im\,ess}(\varphi_C^*\vert\underline{\mathrm{Hom}}_{\bar\delta}(\bar R,C))\subset\underline{\mathrm{Hom}}_\delta(R,C)
\]

b) si $(-)\in\sigma_2$, alors on a l'inclusion opposée $\supset$.

a) $\Delta_\sigma$ stable par équivalence de catégories

b) Une intersection d'\deleted{catégories}\inserted{ensembles} de la forme $\Delta_\sigma$ est de la \unsure{même|mme} forme.

\textcircled{3} $\Delta_{\sigma,\rho}$ est la \deleted{\unsure{$2$}} sous-$2$-catégorie de $\mathrm{Cat}$ dont les objets sont les éléments de $\Delta_\sigma$, et les foncteurs entre \inserted{ses} objets, les foncteurs $\in\deleted{\Delta}\Sigma_\rho$.

a) $\mathrm{Ob}\,\Delta_{\sigma,\rho}$ stable par équivalence

b) Si $C,C'\in\mathrm{Ob}\,\Delta_{\sigma\rho}$, alors $\underline{\mathrm{Hom}}_{\Delta_{\sigma,\rho}}(C,C')\subset\underline{\mathrm{Hom}}_{\mathrm{Cat}}(C,C')$ est une sous-catégorie str. pleine \deleted{de $C$}

(\,La composition des foncteurs de $\Delta_{\sigma\rho}$ induite par celle de $\mathrm{Cat}$)

c) Toute intersection de catégories de la forme $\Delta_{\sigma,\rho}$ est du $=$ type.

% slanted
\margin{sous-2-catégories $\Delta$ ``distinguées''\\ Si a) les $(R,\delta)$ et $(\bar R,\bar\delta)$ \inserted{provenant de} $\in\sigma$ sont aussi \deleted{des}\inserted{parmi les} $(\unsure{\zeta|\Theta},\lambda)\in\rho$\\ b) Si $(\zeta,\lambda)\in\rho$, alors les $\lambda$-limites existent dans $\Delta$, i.e. parmi les \deleted{\unsure{triples de $\sigma$}} quintuples de $\sigma_2$, il y a \deleted{\unsure{ceux}} (avec $\delta=\varnothing,\bar R$ un cône sur $R$) qui expriment l'existence des $\lambda$ limites.}

\uline{Remarque} Soient $\Delta\subset\Delta'$ de la forme $\Delta_{\sigma,\rho}$ et $\Delta_{\sigma',\rho'}$. Soit $R\xrightarrow{\varphi}\bar R$ dans $\Delta'$, et considérons, \inserted{pour $C\in\mathrm{Ob}\,\Delta$,} $u\longmapsto u\circ\varphi$
\[
\underline{\mathrm{Hom}}_{\Delta'}(\bar R,C)\xrightarrow{\varphi^*_{\Delta',C}=\Psi}\underline{\mathrm{Hom}}_{\Delta'}(R,C)
\]

Considérons l'une des conditions suivantes (sur $C$ variable dans $\mathrm{Ob}\,\Delta$)
\[
\left.\begin{aligned}
&\left.\begin{aligned}a)&\ \psi\text{ est fidèle}\\ b)&\ \psi\text{ est surjectif sur les Hom}\end{aligned}\right\}
\text{ conjonction : }\psi\text{ pl. fidèle}\\
&c)\quad\psi\text{ est ess. surjectif sur objets}
\end{aligned}\right\}\text{ conjonction : }\psi\text{ est une équivalence.}
\]

Soit $\Delta_0\subset\Delta$ la sous $2$-catégorie $1$-pleine définie par l'une déterminée des trois conditions précédentes. Alors \uline{$\Delta_0$ est\deleted{elle} distingué ?}

Soit \deleted{\unsure{$\varphi$|$\Delta$}} $\delta$ $(\bar\delta)$ l'ens des \deleted{diagrammes}\inserted{cônes exacts} de $R$ $(\bar R)$ qui sont image des \deleted{$\Theta\to R$} cônes $\lambda$ \deleted{\unsure{ex}} de $\Theta$ par $\Theta\to R$ $(\Theta-\bar R)$ $\lambda$ exacts \deleted{\unsure{dépendant des conditions de|distinctes des conditions dites}} $(\Theta,\lambda)\in\rho'$. Donc $\underline{\mathrm{Hom}}_{\Delta'}(\bar R,C)=\underline{\mathrm{Hom}}_{\bar\delta}(\bar R,C)$, $\underline{\mathrm{Hom}}_{\Delta'}(R,C)=\underline{\mathrm{Hom}}_\delta(R,C)$

% archive sheet number: 13
